\documentclass[UTF8]{ctexart}
\usepackage[a4paper,margin=2.5cm]{geometry}
\usepackage{amsmath,amssymb}
\usepackage{graphicx}
\usepackage{booktabs}
\usepackage{hyperref}

\title{第六题：MNIST 分类中 SGD 与自然梯度训练的对比}
\author{}
\date{\today}

\begin{document}
\emergencystretch=2em
\maketitle

\section{题目与实验安排}

第六题要求构建神经网络对 MNIST 数据集中的 $0$--$9$ 手写数字进行分类，并分别使用随机梯度下降
（SGD）和自然梯度（Natural Gradient, NG）训练，最后比较两种方法的训练效果。程序文件为
\texttt{mnist\_sgd\_ng.py}，它优先读取 \texttt{data/mnist} 中的标准 MNIST IDX 文件，并输出每轮
训练指标、准确率曲线和混淆矩阵。

为了同时展示基础配置与增强配置的差异，本报告综合两组实验：
\begin{center}
\begin{tabular}{lccccc}
\toprule
实验 & 训练样本 & 测试样本 & 隐藏层宽度 & 训练轮数 & 结果目录 \\
\midrule
快速子集实验 & 4000 & 1000 & 64 & 12 & \texttt{results\_mnist} \\
完整增强实验 & 60000 & 10000 & 128 & 30 & \texttt{results\_mnist\_full\_30ep} \\
\bottomrule
\end{tabular}
\end{center}

第一组实验用于观察较小数据规模和较短训练时两种方法的基本表现；第二组实验使用完整 MNIST、较宽隐藏层
和更多训练轮数，用于验证增加训练数据和训练轮数后准确率是否继续提升。

\section{神经网络模型}

两组实验均采用单隐藏层多层感知机（MLP）：
\[
h=\mathrm{ReLU}(xW_1+b_1),\qquad
p_\theta(y\mid x)=\mathrm{softmax}(hW_2+b_2),
\]
其中输入 $x\in\mathbb{R}^{784}$，输出层对应 $10$ 个数字类别。损失函数为带 $L_2$ 正则项的交叉熵：
\[
\mathcal{L}(\theta)=
-\frac{1}{n}\sum_{i=1}^{n}\log p_\theta(y_i\mid x_i)
\;+\;\frac{\lambda}{2}\left(\|W_1\|_2^2+\|W_2\|_2^2\right).
\]

\section{SGD 与自然梯度方法}

SGD 直接沿小批量梯度的反方向更新参数：
\[
\theta_{t+1}=\theta_t-\eta\,\hat g_t,
\]
其中 $\hat g_t=\nabla_\theta \mathcal{L}_B(\theta_t)$ 为当前 minibatch 的梯度。

自然梯度使用模型分布诱导出的 Fisher 信息矩阵 $F(\theta)$ 衡量参数空间的局部几何：
\[
\theta_{t+1}=\theta_t-\alpha\,F(\theta_t)^{-1}\hat g_t.
\]
由于完整 Fisher 矩阵的构造和求逆代价较高，代码采用对角经验 Fisher 近似：
\[
\hat F_t=\beta \hat F_{t-1}+(1-\beta)(\hat g_t\odot \hat g_t),
\]
\[
\theta_{t+1}
=\theta_t-\alpha\left(\mathrm{diag}(\hat F_t)+\lambda_{\mathrm{NG}}I\right)^{-1}\hat g_t.
\]
其中 $\lambda_{\mathrm{NG}}$ 为阻尼项，用于防止 Fisher 对角元素过小导致更新过大。该近似保留了自然梯度
按参数敏感度自适应缩放梯度的思想，同时避免了完整 Fisher 矩阵的存储和求逆。

\section{快速子集实验：4000/1000 样本，12 轮训练}

本实验使用 $4000$ 个训练样本、$1000$ 个测试样本，隐藏层宽度为 $64$，训练 $12$ 轮。运行结果如下：

\begin{center}
\begin{tabular}{lccccc}
\toprule
优化器 & 最终训练准确率 & 最终测试准确率 & 最终测试损失 & 最佳测试准确率 & 最佳轮数 \\
\midrule
SGD & 0.9473 & 0.9070 & 0.3172 & 0.9070 & 12 \\
NG & 0.9150 & 0.8820 & 0.3855 & 0.8820 & 12 \\
\bottomrule
\end{tabular}
\end{center}

\begin{figure}[htbp]
\centering

\includegraphics[width=0.72\linewidth]{results_mnist/sample_digits.png}
\caption{快速子集实验中的 MNIST 数字样例}
\end{figure}

图 1 给出了快速实验中抽取的 MNIST 样例。虽然样本来自真实 MNIST，但训练集只取了 $4000$ 个样本，
类别内书写风格覆盖不足，因此模型泛化能力会受到一定限制。

\begin{figure}[htbp]
\centering
\includegraphics[width=0.86\linewidth]{results_mnist/accuracy_curve.png}
\caption{快速子集实验中 SGD 与 NG 的测试准确率曲线}
\end{figure}

图 2 显示两种方法在 $12$ 轮内都持续上升，并且最佳点都出现在最后一轮，说明此时训练尚未完全收敛。
SGD 的测试准确率达到 $0.9070$，高于 NG 的 $0.8820$。在小样本和较短训练条件下，SGD 的直接梯度更新
更快地提高了训练集拟合程度；NG 由于对角 Fisher 近似和阻尼项会缩放更新步长，前期提升相对保守。

\begin{figure}[htbp]
\centering
\begin{minipage}{0.48\linewidth}
\centering
\includegraphics[width=\linewidth]{results_mnist/confusion_sgd.png}
\caption{快速子集实验的 SGD 混淆矩阵}
\end{minipage}
\hfill
\begin{minipage}{0.48\linewidth}
\centering
\includegraphics[width=\linewidth]{results_mnist/confusion_ng.png}
\caption{快速子集实验的 NG 混淆矩阵}
\end{minipage}
\end{figure}

图 3 和图 4 为快速子集实验的混淆矩阵。两种方法的主要错误集中在形状相近的数字之间，例如 $4$ 与 $9$、
$3$ 与 $5$、$7$ 与 $9$ 等。SGD 的对角线更集中，非对角线误分数量更少，说明在该配置下 SGD 学到的
类别边界更清晰；NG 的错误分布更分散，反映出 $12$ 轮训练和 $4000$ 个训练样本不足以充分发挥其后期
自适应优化优势。

\section{完整增强实验：60000/10000 样本，30 轮训练}

本实验使用全部 $60000$ 个训练样本和 $10000$ 个测试样本，隐藏层宽度增至 $128$，训练轮数增至 $30$。
为使后期训练更稳定，SGD 学习率设为 $0.08$，NG 学习率设为 $0.03$，NG 阻尼设为 $0.5$。运行结果如下：

\begin{center}
\begin{tabular}{lccccc}
\toprule
优化器 & 最终训练准确率 & 最终测试准确率 & 最终测试损失 & 最佳测试准确率 & 最佳轮数 \\
\midrule
SGD & 0.9902 & 0.9771 & 0.0966 & 0.9774 & 29 \\
NG & 0.9863 & 0.9760 & 0.1028 & 0.9760 & 30 \\
\bottomrule
\end{tabular}
\end{center}

\begin{figure}[htbp]
\centering

\includegraphics[width=0.72\linewidth]{results_mnist_full_30ep/sample_digits.png}
\caption{完整增强实验中的 MNIST 数字样例}
\end{figure}

图 5 展示完整 MNIST 实验中的数字样例。相比快速子集实验，完整训练集覆盖了更多书写风格、倾斜程度和局部
形变，这有助于模型学习更稳健的分类边界。

\begin{figure}[htbp]
\centering
\includegraphics[width=0.86\linewidth]{results_mnist_full_30ep/accuracy_curve.png}
\caption{完整增强实验中 SGD 与 NG 的测试准确率曲线}
\end{figure}

图 6 表明两种方法在完整数据上均取得了明显更高的准确率。SGD 最佳测试准确率为 $0.9774$，出现在第
$29$ 轮；NG 最佳测试准确率为 $0.9760$，出现在第 $30$ 轮。相比 $20$ 轮实验，继续训练到 $30$ 轮后，
SGD 的最终测试准确率由 $0.9741$ 提升到 $0.9771$，NG 由 $0.9717$ 提升到 $0.9760$。这说明在完整
MNIST 上，两种方法在 $20$ 轮时尚未完全达到平台期，增加轮数仍能带来小幅但稳定的提升。

\begin{figure}[htbp]
\centering
\begin{minipage}{0.48\linewidth}
\centering
\includegraphics[width=\linewidth]{results_mnist_full_30ep/confusion_sgd.png}
\caption{完整增强实验的 SGD 混淆矩阵}
\end{minipage}
\hfill
\begin{minipage}{0.48\linewidth}
\centering
\includegraphics[width=\linewidth]{results_mnist_full_30ep/confusion_ng.png}
\caption{完整增强实验的 NG 混淆矩阵}
\end{minipage}
\end{figure}

图 7 和图 8 为完整增强实验的混淆矩阵。与快速子集实验相比，两个矩阵的对角线更加集中，非对角线误分
显著减少。SGD 仍略高于 NG，但二者差距已经从快速实验中的 $2.50$ 个百分点缩小到约 $0.11$ 个百分点。
这说明在训练数据更充分、模型容量更大且训练轮数更多时，NG 的对角 Fisher 缩放也能学习到接近 SGD 的
主要类别边界。

\section{两组实验的综合对比}

\begin{center}
\begin{tabular}{llccc}
\toprule
实验配置 & 优化器 & 最终训练准确率 & 最终测试准确率 & 最终测试损失 \\
\midrule
快速配置 & SGD & 0.9473 & 0.9070 & 0.3172 \\
快速配置 & NG & 0.9150 & 0.8820 & 0.3855 \\
完整配置 & SGD & 0.9902 & 0.9771 & 0.0966 \\
完整配置 & NG & 0.9863 & 0.9760 & 0.1028 \\
\bottomrule
\end{tabular}
\end{center}

从表中可以看到，提升准确率的关键不只是增加训练轮数，还包括使用完整训练集和扩大隐藏层宽度。快速子集
实验中，训练样本少、隐藏层较窄、训练轮数较少，SGD 和 NG 的最终测试准确率分别为 $90.70\%$ 与
$88.20\%$。完整增强实验中，二者分别提升到 $97.71\%$ 与 $97.60\%$。因此，更充分的数据覆盖和更大的
模型容量显著降低了泛化误差。

两种优化方法的相对表现也发生了变化。在快速子集实验中，SGD 明显优于 NG，说明在训练时间较短时，普通
SGD 的有效步长更直接，能更快提高拟合程度；NG 的阻尼和对角 Fisher 缩放使其更新更谨慎。完整增强实验中，
NG 与 SGD 的差距大幅缩小，说明当训练轮数增加后，NG 的自适应尺度调整开始发挥作用。不过当前实现只使用
对角 Fisher 近似，忽略了参数之间的相关性，因此仍无法完全体现精确自然梯度的优势。若进一步使用卷积神经
网络、学习率衰减、数据增强或 K-FAC 等更强的 Fisher 近似，准确率仍有继续提升空间。

\section{文件说明}

整理后的第六题文件夹中保留如下主要内容：
\begin{itemize}
  \item \texttt{mnist\_sgd\_ng.py}：主程序，包含 MNIST 读取、MLP、SGD/NG 训练和结果输出；
  \item \texttt{requirements.txt}：运行依赖；
  \item \texttt{data/mnist/}：标准 MNIST IDX 数据文件；
  \item \texttt{results\_mnist/}：快速子集实验的结果图、混淆矩阵和指标文件；
  \item \texttt{results\_mnist\_full\_30ep/}：完整增强实验的结果图、混淆矩阵和指标文件；
  \item \texttt{report.tex} 与 \texttt{report.pdf}：最终报告源文件和编译后的 PDF。
\end{itemize}

\end{document}
